\documentclass[a4paper,zihao=-4]{ctexart}

% ================= 导言区 =================
% 必要的数学宏包
\usepackage{amsmath}   % 核心数学公式
\usepackage{amsthm}    % 定理环境
\usepackage{unicode-math} % 提供 \setmathfont，必须在 ctex 之后加载
\usepackage{libertinus-otf}
\usepackage{graphicx}
\usepackage{algorithm}
\usepackage{algpseudocode}
\usepackage{array,booktabs}
\usepackage{enumitem}
\usepackage{marginnote}
\usepackage{xcolor}
% 引入 tcolorbox 宏包及其皮肤库
\usepackage[skins, breakable]{tcolorbox}
\usepackage{float}
\usepackage{natbib}
\usepackage{hyperref}

\hypersetup{colorlinks=true, linkcolor=blue}
\setcitestyle{authoryear,round,comma}

\renewcommand{\mathbf}[1]{\symbf{#1}}
\renewcommand{\boldsymbol}[1]{\symbfit{#1}}

% 中文正文（需提前安装思源宋体；也可替换为系统自带「宋体」）
\setCJKmainfont{Source Han Serif SC}

% 定义一个名为 "questionbox" 的环境
\newtcolorbox{questionbox}[1][]{ % [1][] 表示有一个可选参数，默认值为空
    enhanced,           % 开启高级渲染
    breakable,          % 允许跨页
    frame hidden,       % 隐藏四周默认边框
    boxrule=0pt,        % 边框粗细归零
    colback=blue!5,     % 背景色：5% 浓度的蓝色
    borderline west={4pt}{0pt}{blue!70}, % 核心：左侧加一条 4pt 粗的深蓝色竖线
    sharp corners,      % 直角（如果喜欢圆角可以去掉这行）
    left=10pt,          % 内部左边距
    right=10pt,         % 内部右边距
    top=8pt,            % 内部上边距
    bottom=8pt,         % 内部下边距
    fontupper=\kaishu,  % 框内中文字体设置为楷体（可选）
    title={#1},         % 允许传入标题
    coltitle=blue!80!black, % 标题颜色
    fonttitle=\bfseries,% 标题加粗
    attach title to upper={\par\vspace{2pt}}, % 标题和正文的间距
}

% 自定义一个小命令，让边注字体变小变灰，不喧宾夺主
\newcommand{\sidenote}[1]{\marginnote{\small\color{gray}#1}}

% 补充定义 \Input 和 \Output，修复未定义错误
\algnewcommand\Input{\item[\textbf{Input:}]}
\algnewcommand\Output{\item[\textbf{Output:}]}
\algnewcommand\Init{\item[\textbf{Initialization:}]}

% 标题信息
\title{论文笔记}
\author{唐巍}
\date{2026 年 7 月}

\begin{document}
\maketitle
\tableofcontents

\newpage

\section{AREM}

\subsection{2-A. REM 模型}

假设 $K$ 个发射机和 $M$ 个传感器随机分布在 ROI 中．\emph{传感器的位置视为先验信息}．第 $m$ 个传感器的位置坐标表示为 $\mathbf{s}_m = [x_m^\mathrm{s}, y_m^\mathrm{s}]^\mathrm{T},\ m = 1, \dots, M$，其中上标 $\mathrm{s}$ 表示传感器的位置坐标，以区别于发射机的位置坐标（用上标 $\mathrm{t}$ 表示）和网格的位置坐标（用上标 $\mathrm{g}$ 表示）．第 $m$ 个传感器在指定时间段和带宽内的 RSS 可写为：
\begin{equation}
	\label{eq:rss_measurement}
	z_m = \xi(\mathbf{s}_m) \sum_{k=1}^{K} p_k f(\mathbf{s}_m, \mathbf{u}_k) + \delta_m
\end{equation}
其中 $\mathbf{u}_k = [x_k^\mathrm{t}, y_k^\mathrm{t}]^\mathrm{T}$ 和 $p_k$ 分别表示第 $k$ 个发射机的位置坐标和功率．$\xi(\mathbf{s}_m)$ 表示第 $m$ 个传感器位置处的阴影衰落，服从\emph{对数正态分布}．$\delta_m$ 是表示测量误差的高斯随机变量．路径损耗函数 $f(\mathbf{s}_m, \mathbf{u}_k)$ 由路径损耗模型指定：
\begin{equation}
	\label{eq:path_loss_function}
	f(\mathbf{s}_m, \mathbf{u}_k) =
	\begin{cases}
		1,                              & \text{for } d_{mk} \leq d_0 \\
		\left(d_0 / d_{mk}\right)^\eta, & \text{for } d_{mk} > d_0
	\end{cases}
\end{equation}
其中 $d_{mk} = \|\mathbf{s}_m - \mathbf{u}_k\|_2 = \sqrt{(x_m^\mathrm{s} - x_k^\mathrm{t})^2 + (y_m^\mathrm{s} - y_k^\mathrm{t})^2}$ 是发射器与传感器之间的欧氏距离．$d_0$ 表示天线远场参考距离．$\eta$ 表示路径损耗指数，由环境决定．基于式\eqref{eq:path_loss_function}，路径损耗向量 $\mathbf{a}(\mathbf{u}_k)$ 表示为
\begin{equation}
	\label{eq:path_loss_vector}
	\mathbf{a}(\mathbf{u}_k) = \left[ f(\mathbf{s}_1, \mathbf{u}_k), \dots, f(\mathbf{s}_M, \mathbf{u}_k) \right]^\mathrm{T} \in \mathbb{R}^M
\end{equation}
并且路径损耗字典 $\mathbf{A}(\mathbf{u})$ 表示为
\begin{equation}
	\label{eq:path_loss_dictionary}
	\mathbf{A}(\mathbf{u}) = \left[ \mathbf{a}(\mathbf{u}_1), \dots, \mathbf{a}(\mathbf{u}_K) \right] \in \mathbb{R}^{M \times K}.
\end{equation}
因此，式\eqref{eq:rss_measurement}可以写成矩阵形式如下：
\begin{equation}
	\label{eq:rss_matrix_form}
	\mathbf{z} = \xi(\mathbf{s}) \odot \left( \mathbf{A}(\mathbf{u}) \mathbf{p} \right) + \boldsymbol{\delta}
\end{equation}
其中 $\mathbf{z} = [z_1, \dots, z_M]^\mathrm{T}$，$\boldsymbol{\delta} = [\delta_1, \dots, \delta_M]^\mathrm{T}$，$\xi(\mathbf{s}) = [\xi(\mathbf{s}_1), \dots, \xi(\mathbf{s}_M)]^\mathrm{T}$，且 $\mathbf{p} = [p_1, \dots, p_K]^\mathrm{T}$．式\eqref{eq:rss_matrix_form}可改写如下
\begin{equation}
	\label{eq:rss_rewritten}
	\mathbf{z} = \mathbf{A}(\mathbf{u}) \mathbf{p} + \mathbf{n}
\end{equation}
其中 $\mathbf{n} = \left( \xi(\mathbf{s}) - \mathbf{1}_M \right) \odot \left( \mathbf{A}(\mathbf{u}) \mathbf{p} \right) + \boldsymbol{\delta}$．当对数正态分布 $\xi(\mathbf{s})$ 的方差足够小时，$\mathbf{n}$ 可近似为高斯分布，用以解释阴影衰落效应、测量误差以及建模失配．

\subsection{2-B. 稀疏恢复模型}

基于发射机在感兴趣区域（ROI）中的稀疏空间分布，式\eqref{eq:rss_rewritten}可被建模为一个稀疏恢复问题．具体而言，ROI 通常被离散化为具有相等网格间距 $\Delta$ 的网格．假设所有发射机精确地位于网格上，且网格总数为 $N$．在此假设下，式\eqref{eq:rss_rewritten}可表示为
\begin{equation}
	\label{eq:sparse_grid_model}
	\mathbf{z} = \mathbf{A}(\boldsymbol{\theta}) \boldsymbol{\omega} + \mathbf{n}
\end{equation}
其中 $\mathbf{A}(\boldsymbol{\theta}) = [\mathbf{a}(\boldsymbol{\theta}_1), \dots, \mathbf{a}(\boldsymbol{\theta}_i), \dots, \mathbf{a}(\boldsymbol{\theta}_N)] \in \mathbb{R}^{M \times N}$，且 $\boldsymbol{\theta}_i = [x_i^\mathrm{g}, y_i^\mathrm{g}]^\mathrm{T}$ 表示第 $i$ 个网格的位置坐标．$\boldsymbol{\omega} = [\omega_1, \dots, \omega_i, \dots, \omega_N]^\mathrm{T} \in \mathbb{R}^N$ 满足：
\begin{equation}
	\label{eq:sparse_vector_def}
	\omega_i =
	\begin{cases}
		p_k, & \text{第 } k \text{ 个发射机在第 } i \text{ 个网格上} \\
		0,   & \text{第 } i \text{ 个网格上没有发射机}.
	\end{cases}
\end{equation}
由于 ROI 中的发射机是稀疏分布的，$\boldsymbol{\omega}$ 是一个稀疏度为 $K$ 的稀疏向量（$K$ 表示 $\boldsymbol{\omega}$ 中非零元素的个数，且 $K \ll N$），相应的稀疏恢复模型可表示如下：
\begin{equation}
	\label{eq:on_grid_problem}
	\begin{aligned}
		\min_{\boldsymbol{\omega}} \quad & \| \mathbf{z} - \mathbf{A} \boldsymbol{\omega} \|_2^2 \\
		\text{s.t.} \quad                & \| \boldsymbol{\omega} \|_0 \leq K.
	\end{aligned}
\end{equation}
式\eqref{eq:on_grid_problem}是一个网格（on-grid）模型，仅需求解 $\boldsymbol{\omega}$．在实际场景中，发射机可位于 ROI 内的任意位置，这使得预定义的离散网格难以包含所有潜在的发射机位置．离网效应会导致网格模型出现估计偏差．该离网问题（off-grid）可被建模为一个字典学习问题
\begin{equation}
	\label{eq:off_grid_problem}
	\begin{aligned}
		\min_{\mathbf{A}(\boldsymbol{\theta}), \boldsymbol{\omega}} \quad & \| \mathbf{z} - \mathbf{A}(\boldsymbol{\theta}) \boldsymbol{\omega} \|_2^2 \\
		\text{s.t.} \quad                                                 & \| \boldsymbol{\omega} \|_0 \leq K.
	\end{aligned}
\end{equation}
需要指出的是，与式\eqref{eq:on_grid_problem}相比，式\eqref{eq:off_grid_problem}是一个非凸的双线性逆问题，其中字典矩阵 $\mathbf{A}(\boldsymbol{\theta})$ 和向量 $\boldsymbol{\omega}$ 都需要求解．具体地，$\boldsymbol{\theta}$ 在连续空间中进行估计．显然，求解式\eqref{eq:off_grid_problem}有助于提高 REM 重建的精度．在第三部分中，提出了 AREM 来求解式\eqref{eq:off_grid_problem}．
%\begin{questionbox}[关键问题 (Question)]
%    为什么在 Transformer 架构中，Layer Normalization 通常放置在多头注意力机制之前（Pre-LN），而不是像原始论文那样放置在之后（Post-LN）？这两种方式对梯度回传有什么影响？
%\end{questionbox}
%在经典的凸优化理论中，我们经常需要判断矩阵的正定性 \sidenote{这是一个重要的边注：正定矩阵的所有特征值必须大于零．}．如果 Hessian 矩阵是半正定的，那么目标函数就是一个凸函数．

\subsection{3-A. 方法描述}

AREM 基于交替迭代更新的思想求解非凸双线性逆问题 \eqref{eq:off_grid_problem}．迭代过程可概括如下：在第 $k$ 次迭代开始时，根据与残差信号的相关性获得对一个新的发射机位置参数的粗略估计，并将其加入支持集 $\mathcal{T}_k = \{ \boldsymbol{\theta}_1, \ldots, \boldsymbol{\theta}_i, \ldots, \boldsymbol{\theta}_{|\mathcal{T}_k|} \}$．当得到 $\mathcal{T}_k$ 时，即 $\mathbf{A}_{\mathcal{T}_k}(\boldsymbol{\theta}) = [\mathbf{a}(\boldsymbol{\theta}_1), \dots, \mathbf{a}(\boldsymbol{\theta}_i), \dots, \mathbf{a}(\boldsymbol{\theta}_{|\mathcal{T}_k|})] \in \mathbb{R}^{M \times |\mathcal{T}_k|}$ 已获得（在下文中，$\mathbf{A}_{\mathcal{T}_k}(\boldsymbol{\theta})$ 简写为 $\mathbf{A}_{\mathcal{T}_k}$），$\boldsymbol{\omega}_k$ 可通过最小二乘估计进行更新．当 $\boldsymbol{\omega}_k$ 被估计后，$\mathbf{A}_{\mathcal{T}_k}$ 可通过连续空间中的局部加细进行更新．然后，为保持稀疏性，通过剪枝移除不满意的估计．最后，更新残差 $\boldsymbol{\gamma}_k$．详细过程描述如下．

\paragraph{1. 粗略估计}
假设 ROI 被离散化为网格，网格总数为 $N$．在第 $k$ 次迭代开始时，通过选择与残差 $\boldsymbol{\gamma}_{k-1}$ 相关性最大的 $\mathbf{a}(\boldsymbol{\theta}_i)$，可以获得新发射机位置参数的粗略估计．由于 $\mathbf{A}(\boldsymbol{\theta})$ 列间存在强相关性（如图 2(a) 所示），需进行以下数据预处理以降低相关性，从而提高粗略估计的准确性：
\begin{equation}
    \label{eq:preprocessed_model}
    \mathbf{T}\mathbf{z} = \mathbf{T}\mathbf{A}(\boldsymbol{\theta})\boldsymbol{\omega} + \mathbf{T}\mathbf{n}
\end{equation}
其中 $\mathbf{T} = \mathbf{Q}\mathbf{A}^{\dagger}(\boldsymbol{\theta})$，$\mathbf{Q}$ 是 $\mathbf{A}(\boldsymbol{\theta})$ 的正交基．令 $\mathbf{A}'(\boldsymbol{\theta}) = \mathbf{T}\mathbf{A}(\boldsymbol{\theta})$ 且 $\boldsymbol{\gamma}_{k-1}' = \mathbf{T}\boldsymbol{\gamma}_{k-1}$．从图 2(b) 中可以观察到，$\mathbf{A}'(\boldsymbol{\theta})$ 列间的相关性显著降低．粗略估计可表示为
\begin{equation}
    \label{eq:rough_estimation}
    \boldsymbol{\theta}_k \leftarrow \arg\max \left| \left\langle \mathbf{a}'(\boldsymbol{\theta}_i), \boldsymbol{\gamma}_{k-1}' \right\rangle \right|, \quad i = 1, \dots, N.
\end{equation}
\paragraph{2. 更新支撑集}
将参数 $\boldsymbol{\theta}_k$ 添加到支撑集 $\mathcal{T}_{k-1}$ 中，形成新的支撑集 $\mathcal{T}_k = \{ \boldsymbol{\theta}_1, \ldots, \boldsymbol{\theta}_i, \ldots, \boldsymbol{\theta}_{|\mathcal{T}_k|} \}$．
\paragraph{3. 交替更新过程}
交替更新过程包含四个步骤．第一步是更新 $\boldsymbol{\omega}_k$．当估计出 $\mathcal{T}_k$，即得到 $\mathbf{A}_{\mathcal{T}_k}(\boldsymbol{\theta})$ 后，可通过最小二乘估计求解以下优化问题来更新 $\boldsymbol{\omega}_k$：
\begin{equation}
    \label{eq:omega_update}
    \begin{aligned}
        \boldsymbol{\omega}_k &= \arg\min_{\boldsymbol{\omega}} \left\| \mathbf{z} - \mathbf{A}_{\mathcal{T}_k} \boldsymbol{\omega} \right\|_2^2 = \mathbf{A}_{\mathcal{T}_k}^{\dagger} \mathbf{z} \\
        &= \left( \mathbf{A}_{\mathcal{T}_k}^{\mathrm{T}} \mathbf{A}_{\mathcal{T}_k} \right)^{-1} \mathbf{A}_{\mathcal{T}_k}^{\mathrm{T}} \mathbf{z}.
    \end{aligned}
\end{equation}
第二步是剪枝．为保持稀疏性，$\boldsymbol{\omega}_k$ 中小于足够小正数的元素以及 $\mathcal{T}_k$ 中超出 ROI 范围的元素将从 $\boldsymbol{\omega}_k$ 和 $\mathcal{T}_k$ 中移除．剪枝过程表示为
\begin{equation}
    \label{eq:prune}
    \boldsymbol{\omega}_k, \mathcal{T}_k \leftarrow \operatorname{Prune}(\boldsymbol{\omega}_k, \mathcal{T}_k).
\end{equation}
第三步是更新 $\mathbf{A}_{\mathcal{T}_k}$．当估计出 $\boldsymbol{\omega}_k$ 后，可通过局部细化在连续空间中优化 $\mathbf{A}_{\mathcal{T}_k}$，如下所示：
\begin{equation}
    \label{eq:A_optimization}
    \min_{\mathbf{A}_{\mathcal{T}_k}} \left\| \mathbf{z} - \mathbf{A}_{\mathcal{T}_k} \boldsymbol{\omega}_k \right\|_2^2 = \min_{\mathbf{A}_{\mathcal{T}_k}} \left\| \mathbf{z} - \mathbf{A}_{\mathcal{T}_k} \mathbf{A}_{\mathcal{T}_k}^{\dagger} \mathbf{z} \right\|_2^2.
\end{equation}
式 \eqref{eq:A_optimization} 可通过梯度下降算法求解．定义目标函数 $H(\boldsymbol{\theta})$ 如下：
\begin{equation}
    \label{eq:objective_H}
    \begin{aligned}
        H(\boldsymbol{\theta}) &= \left\| \mathbf{z} - \mathbf{A}_{\mathcal{T}_k} \mathbf{A}_{\mathcal{T}_k}^{\dagger} \mathbf{z} \right\|_2^2 \\
        &= \mathbf{z}^{\mathrm{T}} \mathbf{z} - \mathbf{z}^{\mathrm{T}} \mathbf{A}_{\mathcal{T}_k} \mathbf{A}_{\mathcal{T}_k}^{\dagger} \mathbf{z}.
    \end{aligned}
\end{equation}
$H(\boldsymbol{\theta})$ 的梯度可表示为
\begin{equation}
    \label{eq:gradient_H}
    \nabla H(\boldsymbol{\theta}) = \begin{bmatrix}
        \frac{\partial H(\boldsymbol{\theta})}{\partial x_1^\mathrm{g}}, \frac{\partial H(\boldsymbol{\theta})}{\partial y_1^\mathrm{g}} \\
        \vdots \\
        \frac{\partial H(\boldsymbol{\theta})}{\partial x_i^\mathrm{g}}, \frac{\partial H(\boldsymbol{\theta})}{\partial y_i^\mathrm{g}} \\
        \vdots \\
        \frac{\partial H(\boldsymbol{\theta})}{\partial x_{|\mathcal{T}_k|}^\mathrm{g}}, \frac{\partial H(\boldsymbol{\theta})}{\partial y_{|\mathcal{T}_k|}^\mathrm{g}}
    \end{bmatrix}^{\mathrm{T}}
\end{equation}
其中 $\partial H(\boldsymbol{\theta})/\partial x_i^\mathrm{g}$ 可通过下式计算：
\begin{equation}
    \label{eq:partial_derivative_x}
    \begin{aligned}
        \frac{\partial H(\boldsymbol{\theta})}{\partial x_i^\mathrm{g}} &= \mathbf{z}^{\mathrm{T}} \left( \mathbf{R}_i + \mathbf{R}_i^{\mathrm{T}} \right) \mathbf{z} \\
        \mathbf{R}_i &= \left[ \mathbf{I} - \mathbf{A}_{\mathcal{T}_k} \mathbf{A}_{\mathcal{T}_k}^{\dagger} \right] \frac{\partial \mathbf{A}_{\mathcal{T}_k}}{\partial x_i^\mathrm{g}} \mathbf{A}_{\mathcal{T}_k}^{\dagger} \\
        \frac{\partial \mathbf{A}_{\mathcal{T}_k}}{\partial x_i^\mathrm{g}} &= \left[ \mathbf{0}, \ldots, \frac{\partial \mathbf{a}(\boldsymbol{\theta}_i)}{\partial x_i^\mathrm{g}}, \ldots, \mathbf{0} \right] \\
        \frac{\partial \mathbf{a}(\boldsymbol{\theta}_i)}{\partial x_i^\mathrm{g}} &= \left[ \frac{\partial f(\mathbf{s}_1, \boldsymbol{\theta}_i)}{\partial x_i^\mathrm{g}}, \ldots, \frac{\partial f(\mathbf{s}_M, \boldsymbol{\theta}_i)}{\partial x_i^\mathrm{g}} \right]^{\mathrm{T}} \\
        \frac{\partial f(\mathbf{s}_m, \boldsymbol{\theta}_i)}{\partial x_i^\mathrm{g}} &= \eta d_0^\eta \frac{x_m^\mathrm{s} - x_i^\mathrm{g}}{\| \mathbf{s}_m - \boldsymbol{\theta}_i \|_2^{\eta+2}}.
    \end{aligned}
\end{equation}
$\partial H(\boldsymbol{\theta})/\partial y_i^\mathrm{g}$ 可同理计算．基于 $H(\boldsymbol{\theta})$ 的可微性，$\mathcal{T}_k$ 中的元素通过连续空间中的梯度下降算法更新如下：
\begin{equation}
    \label{eq:gradient_update}
    \begin{bmatrix}
        \boldsymbol{\theta}_1^{\mathrm{T}} \\
        \vdots \\
        \boldsymbol{\theta}_i^{\mathrm{T}} \\
        \vdots \\
        \boldsymbol{\theta}_{|\mathcal{T}_k|}^{\mathrm{T}}
    \end{bmatrix}
    \leftarrow
    \begin{bmatrix}
        \boldsymbol{\theta}_1^{\mathrm{T}} \\
        \vdots \\
        \boldsymbol{\theta}_i^{\mathrm{T}} \\
        \vdots \\
        \boldsymbol{\theta}_{|\mathcal{T}_k|}^{\mathrm{T}}
    \end{bmatrix}
    - \alpha \nabla H^{\mathrm{T}}(\boldsymbol{\theta})
\end{equation}
其中 $\alpha$ 表示梯度下降步长．随后，得到 $\mathbf{A}_{\mathcal{T}_k} = [\mathbf{a}(\boldsymbol{\theta}_1), \dots, \mathbf{a}(\boldsymbol{\theta}_i), \dots, \mathbf{a}(\boldsymbol{\theta}_{|\mathcal{T}_k|})]$ 的更新．在最后一步，更新残差 $\boldsymbol{\gamma}_k = \mathbf{z} - \mathbf{A}_{\mathcal{T}_k}\boldsymbol{\omega}_k$．当残差的变化小于足够小的正阈值（$\mathrm{th}_{\text{error}}$）时，交替更新过程停止，并执行下一次迭代．所提出的 AREM 方法的伪代码总结在算法 \ref{alg:arem} 中．

\begin{algorithm}[htbp]
    \caption{AREM Method}
    \label{alg:arem}
    \begin{algorithmic}[1]
        \Input{$\mathbf{z}$, $\mathbf{A}$, $\mathbf{A}'$, $\mathbf{T}$}
        \Init $\boldsymbol{\gamma}_0 = \mathbf{z}$, $\mathcal{T}_0 = \emptyset$, $\varepsilon_0 = \left\| \boldsymbol{\gamma}_0 \right\|_2$, $k = 1$

        \While{termination conditions not met}
        \State Rough Estimation:
        \State \quad Data Preprocessing: $\boldsymbol{\gamma}_{k-1}' = \mathbf{T} \boldsymbol{\gamma}_{k-1}$
        \State \quad $\boldsymbol{\theta}_k \leftarrow \arg\max \left| \left\langle \mathbf{a}'(\boldsymbol{\theta}_i), \boldsymbol{\gamma}_{k-1}' \right\rangle \right|,\ i = 1, \dots, N$

        \State Update Support: $\mathcal{T}_k \leftarrow \mathcal{T}_{k-1} \cup \{\boldsymbol{\theta}_k\}$

        \State Alternating Update Process:
        \While{$|\varepsilon_{k-1}| \geq \mathrm{th}_{\text{error}}$}
        \State Update $\boldsymbol{\omega}_k$:
        \State \quad $\boldsymbol{\omega}_k = \left( \mathbf{A}_{\mathcal{T}_k}^{\mathrm{T}} \mathbf{A}_{\mathcal{T}_k} \right)^{-1} \mathbf{A}_{\mathcal{T}_k}^{\mathrm{T}} \mathbf{z}$

        \State Pruning:
        \State \quad $\boldsymbol{\omega}_k,\ \mathcal{T}_k \leftarrow \text{Prune}(\boldsymbol{\omega}_k, \mathcal{T}_k)$

        \State Update $\mathbf{A}_{\mathcal{T}_k}$:
        \State \quad $\mathbf{A}_{\mathcal{T}_k} \leftarrow \text{Gradient\_update}(H(\boldsymbol{\theta}))$, by Eq.~\eqref{eq:gradient_H},
        \State \quad Eq.~\eqref{eq:partial_derivative_x}, and Eq.~\eqref{eq:gradient_update}

        \State Update Residual:
        \State \quad $\boldsymbol{\gamma}_k = \mathbf{z} - \mathbf{A}_{\mathcal{T}_k} \boldsymbol{\omega}_k$
        \State \quad $\varepsilon_k = \left\| \boldsymbol{\gamma}_k \right\|_2 - \left\| \boldsymbol{\gamma}_{k-1} \right\|_2$
        \EndWhile

        \State $k = k + 1$
        \EndWhile

        \Output{$\mathcal{T}_k$, $\boldsymbol{\omega}_k$}
    \end{algorithmic}
\end{algorithm}

\subsection{数值验证}
\subsubsection*{3-B. 粗略估计的影响}

AREM 是一种非凸优化方法，$H(\boldsymbol{\theta})$ 的收敛性与初始粗略估计密切相关．因此，本节基于以下场景分析粗略估计对 AREM 收敛性能的影响：假设感兴趣区域大小为 $400 \times 400\,\mathrm{m}$，网格间距为 $20\,\mathrm{m}$（$\Delta = 20\,\mathrm{m}$）．感兴趣区域内分布有 200 个传感器（$M = 200$）．路径损耗指数设为 $2$（$\eta = 2$），天线远场参考距离设为 $5\,\mathrm{m}$（$d_0 = 5\,\mathrm{m}$）．接下来，分单发射器和多发射器两种情况分析粗略估计的影响．\textbf{数值任务}如下：

\begin{enumerate}
    \item 假设单个发射器的实际位置坐标为 $[191, 191]^\mathrm{T}$，该坐标不在预定义网格（分辨率为 $20\,\mathrm{m}$）上．当 $\Delta = 20\,\mathrm{m}$ 时，粗略估计结果为 $\boldsymbol{\theta} = ?$．然后令 $F(\boldsymbol{\theta}) = -H(\boldsymbol{\theta})$，执行交替更新，验证 $F(\boldsymbol{\theta})$ 是否在发射器实际位置处达到最大值，即验证局部细化后 $\boldsymbol{\theta} = [191, 191]^\mathrm{T}$ 是否成立？然后，可视化 $F(\boldsymbol{\theta})$ 的 3D 视图．
    \item 假设发射机数量 $K$ 从 3 变化到 6，并随机分布在 ROI 中．比较两种策略下 AREM 的定位精度：策略 1 在粗估计步骤中将网格密度加倍，而策略 2 则不进行加倍．具体来说，随机生成分辨率 $1\,\mathrm{m}$ 的 $K$ 个发射机坐标，在 $K$ 从 3 到 6 的四种情形下，比较策略 1 和策略 2 估计的发射机坐标 $\boldsymbol{\theta}$ 与真实坐标的差异．实验结果是否能反映：随着发射机数量的增加，，策略 1 仍能准确估计每个发射机的位置，而策略 2 下的位置估计出现不准确？
\end{enumerate}

%\begin{enumerate}
%    \item \textbf{单发射器情况：}
%    假设发射器的实际位置坐标为 $[191, 191]^\mathrm{T}$，该坐标不在预定义网格上．如图 3(a) 所示，当 $\Delta = 20\,\mathrm{m}$ 时，粗略估计结果为 $\boldsymbol{\theta} = [200, 200]^\mathrm{T}$，该结果位于接近发射器实际位置的网格点上．为便于观察，令 $F(\boldsymbol{\theta}) = -H(\boldsymbol{\theta})$．如图 4 所示，$F(\boldsymbol{\theta})$ 在发射器实际位置处达到最大值，即 $\boldsymbol{\theta} = [191, 191]^\mathrm{T}$（如图 4(b) 中\textbf{红色}实心圆点所示）．
%    如图 4(b) 所示，当 $\boldsymbol{\theta}$ 在红色虚线框所围区域内变化时，$F(\boldsymbol{\theta})$ 表现出单峰函数特性．在此条件下，通过交替更新过程中的梯度下降算法，$F(\boldsymbol{\theta})$ 能够收敛到全局最优．可以清楚地看到，粗略估计结果（如图 4(b) 中黑色实心圆点所示）位于红色虚线框内，这证明了 AREM 在单发射器情况下具有良好的全局收敛能力．

%    \item \textbf{多发射器情况：}
%    从图 4(b) 可以观察到，当 $\boldsymbol{\theta}$ 在红色虚线框内的蓝色区域变化时，由于天线远场参考距离内的路径损耗函数被假定为常数，$F(\boldsymbol{\theta})$ 的值变化非常小，导致梯度 $\nabla F(\boldsymbol{\theta})$ 趋近于零向量．在多发射器情况下，如果 $\boldsymbol{\theta}$ 的粗略估计位于蓝色区域，$\nabla F(\boldsymbol{\theta})$ 容易受到来自其他发射器旁瓣的干扰．这可能会阻止 $F(\boldsymbol{\theta})$ 收敛到全局最优或降低其收敛速度．
%    然而，如果粗略估计结果位于 $\nabla F(\boldsymbol{\theta})$ 的主瓣区域（如图 4(b) 中\textbf{红色实线框}所示），$\nabla F(\boldsymbol{\theta})$ 的一致性和稳定性会显著增强，这对促进 $F(\boldsymbol{\theta})$ 的收敛起着非常重要的作用．如图 3(b) 所示，当网格密度加倍时（即 $\Delta$ 从 $20\,\mathrm{m}$ 减小到 $10\,\mathrm{m}$），粗略估计结果变为 $\boldsymbol{\theta} = [190, 190]^\mathrm{T}$（如图 4(b) 中\textbf{绿色}实心圆点所示），该结果位于主瓣区域．因此，在粗略估计步骤中增加网格密度可以提高估计精度，从而增强 AREM 的收敛性．
%\end{enumerate}

%为进一步验证所提策略的有效性，比较了两种策略下 AREM 的收敛性：策略 1 在粗估计步骤中将网格密度加倍，而策略 2 则不进行加倍．如图 5 所示，虽然当 $K = 3$ 时，两种策略均能准确估计发射机位置，但随着发射机数量的增加，旁瓣的影响导致策略 2 下的位置估计出现不准确．相比之下，策略 1 下仍能准确估计每个发射机的位置，证明所提策略能显著提高 AREM 的收敛性．因此，后续章节中 AREM 采用策略 1．


\subsubsection*{3-C. 参数估计性能}

为评估参数估计性能，我们首先定义位置估计误差、功率估计误差以及正确计数发射机的概率如下：
\begin{equation}
    \label{eq:position_error}
    \text{position\_error} = \frac{1}{K} \sum_{k=1}^{K} \left\| \dot{\mathbf{u}}_k - \mathbf{u}_k \right\|_2
\end{equation}
其中 $\dot{\mathbf{u}}_k$ 和 $\mathbf{u}_k$ 分别表示估计位置和实际位置，$K$ 为发射机数量．
\begin{equation}
    \label{eq:power_error}
    \text{power\_error} = \frac{1}{K} \sum_{k=1}^{K} \left| \frac{\dot{p}_k - p_k}{p_k} \right|^2
\end{equation}
其中 $\dot{p}_k$ 和 $p_k$ 分别表示估计功率和实际功率．
\begin{equation}
    \label{eq:success_probability}
    P_r = \frac{T_{\mathrm{success}}}{T}
\end{equation}
其中 $T$ 表示独立实验的总次数（本文中 $T=200$），$T_{\mathrm{success}}$ 是发射机数量被正确估计的次数．位置误差和功率误差是基于发射机正确计数来计算的．

在上述参数设置的基础上，假设发射机数量 $K$ 从 2 变化到 9，每个发射机功率为 $0\,\mathrm{dBm}$，并随机分布在 ROI 中．

\textbf{数值任务}：计算 AREM 的位置估计误差、功率估计误差以及正确计数发射机的概率随发射机数量的变化情况．


%然后，使用 AREM、VBEM 和正交匹配追踪（OMP）估计发射机参数．

%考虑到传感器分布对 VBEM 性能的影响，我们还进行了一项仿真，评估 VBEM 在均匀传感器分布下的估计性能．结果如图 6 所示，其中 VBEM-Random 和 VBEM-Uniform 分别表示随机和均匀传感器分布下的估计结果．可以看出，当传感器均匀分布时，VBEM 的性能显著提升．尽管如此，AREM 仍然表现出最佳性能．与 VBEM-Uniform 和 OMP 相比，AREM 的位置估计误差分别降低了约 27\% 和 93\%，功率估计误差分别降低了约 6\% 和 99\%．此外，正确计数概率分别提高了约 6\% 和 11\%．


\subsubsection*{4-A. 仿真与讨论}
%本节通过数值仿真和实际场景实验验证 AREM 的 REM 重建性能，并将其与现有方法进行比较．
本节通过数值仿真验证 AREM 的 REM 重建性能．为了评估 REM 重建精度，定义平均加权误差（AWErr）如下：
\begin{equation}
    \label{eq:awerr}
    \mathrm{AWErr} = \sum_{i=1}^{N} \lambda_i \left| P_i - \dot{P}_i \right|
\end{equation}
其中 $P_i$ 是第 $i$ 个网格的实际 RSS（单位 $\mathrm{dBm}$），$\dot{P}_i$ 是该网格的估计 RSS．权重 $\lambda_i$ 量化每个网格上误差的重要性，$\lambda_i = 1/N$ 表示平均绝对误差．因此，如果第 $i$ 个网格的实际 RSS $P_i$ 较大，则该网格对应的误差被赋予更高的权重．
在这种情况下，假设权重 $\lambda_i$ 与 $P_i$ 成正比，并满足约束条件 $\sum_{i=1}^{N} \lambda_i = 1$．基于此推理，我们设定$\lambda_i = P_i \big/ \sum_{i=1}^{N} P_i$．

本节通过数值仿真评估参数对 AREM 重建性能的影响．假设 ROI 大小为 $400 \times 400\,\mathrm{m}$，网格间距为 $5\,\mathrm{m}$（$\Delta = 5\,\mathrm{m}$）．三个发射器（$K = 3$）随机分布在 ROI 中，每个发射功率为 $0\,\mathrm{dBm}$．200 个传感器随机分布在 ROI 中．路径损耗指数设为 $2$（$\eta = 2$），天线远场参考距离设为 $1\,\mathrm{m}$（$d_0 = 1\,\mathrm{m}$）．本节进行 200 次独立蒙特卡洛仿真，并使用平均 AWErr 作为评估 REM 重建性能的指标．\textbf{数值任务}如下：

%本节通过数值仿真评估参数对 AREM 重建性能的影响，并与 TPS、OMP 和 VBEM 进行比较．假设 ROI 大小为 $400 \times 400\,\mathrm{m}$，网格间距为 $5\,\mathrm{m}$（$\Delta = 5\,\mathrm{m}$）．三个发射器（$K = 3$）随机分布在 ROI 中，每个发射功率为 $0\,\mathrm{dBm}$．两百个传感器随机分布在 ROI 中．路径损耗指数设为 $2$（$\eta = 2$），天线远场参考距离设为 $1\,\mathrm{m}$（$d_0 = 1\,\mathrm{m}$）．在以下实验中，参数的变化将根据需要说明．本节进行 200 次独立仿真，并使用平均 AWErr 作为评估 REM 重建性能的指标．

\begin{enumerate}
    \item \textbf{传感器数量的影响：}
    首先研究重建性能与传感器数量之间的关系．传感器数量从 50 增加到 400（以50为增量），计算对应的 AWErr．

%    如图 9 所示，AREM 的重建误差最低，并随着传感器数量的增加迅速趋近于 $0\,\mathrm{dBm}$．虽然 VBEM 的重建误差低于 TPS 和 OMP，但当传感器数量超过 250 时，其误差稳定在约 $2\,\mathrm{dBm}$ 且不再下降．由于离网效应，OMP 表现出最差的重建性能．

    \item \textbf{发射器数量的影响：}
    接下来研究重建性能与发射器数量之间的关系．当发射器数量从 4 变化到 10 时，计算对应的 AWErr．

%    如图 10 所示，当发射器数量从 4 变化到 10 时，AREM 性能最佳，其次是 VBEM、TPS 和 OMP．随着发射器数量的增加，VBEM 的重建误差逐渐上升，而 AREM 的重建误差始终保持较低水平．

    \item \textbf{SNR 的影响：}
如前文所述，噪声 $\mathbf{n}$ 表示阴影衰落效应、测量误差和建模失配．这里研究重建性能与 SNR 之间的关系．当 SNR 从 $5\,\mathrm{dB}$ 变化到 $50\,\mathrm{dB}$ （以5为增量）时，计算对应的 AWErr．其中，分别计算传感器数量为 100,200,300 的三种情况．
%    如第 2-B 节所述，噪声 $\mathbf{n}$ 表示阴影衰落效应、测量误差和建模失配．这里研究重建性能与 SNR 之间的关系．如图 11 所示，与 OMP、TPS 和 VBEM 相比，当 SNR 从 $5\,\mathrm{dB}$ 变化到 $50\,\mathrm{dB}$ 时，AREM 的重建误差分别降低了约 93\%、79\% 和 33\%．随着传感器数量的增加，AREM 的重建误差进一步降低，表明其对噪声干扰具有良好的鲁棒性．

    \item \textbf{发射功率的影响：}
    接着研究重建性能与发射功率之间的关系．当发射功率分别在 [0,5]、[0,10]、[0,15]、[0,20] 内生成时，计算对应的 AWErr．其中在每次独立实验中，让 3 个发射器的功率 $p_k$ 在该区间内服从均匀分布．例如，对于[0,10]情形，某次实验的功率可能是 {1.2 dBm, 18.5 dBm, 9.0 dBm}．
%    如图 12 所示，与 OMP、TPS 和 VBEM 相比，AREM 表现出优越的性能，因为它通过在每次迭代中更新残差，有效减少了高功率信号与低功率信号之间的干扰．

    \item \textbf{ROI 的影响：}
    接下来，研究 REM 重建性能与 ROI 之间的关系．假设传感器数量为 500（$M = 500$），且它们随机分布在 ROI 中．
    当 ROI 从 $1000 \times 1000\,\mathrm{m}$ 变化到 $10000 \times 10000\,\mathrm{m}$ （以1000为增量）时，计算对应的 AWErr．
%    如图 13 所示，当 ROI 从 $1000 \times 1000\,\mathrm{m}$ 变化到 $10000 \times 10000\,\mathrm{m}$ 时，AREM 的重建误差仍然最低．值得注意的是，随着 ROI 从 $4000 \times 4000\,\mathrm{m}$ 增加到 $10000 \times 10000\,\mathrm{m}$，AREM 的重建误差保持不变．

    \item \textbf{路径损耗指数的影响：}
    路径损耗指数 $\eta$ 反映了传播环境的特性．例如，在建筑物内的视距传播条件下，路径损耗指数通常在 1.6 到 1.8 之间；在自由空间中约为 2；而在城市蜂窝阴影条件下，通常超过 3．为了研究重建性能与路径损耗指数之间的关系，计算当 $\eta$ 从 1.5 到 4 （以0.5为增量）时对应的 AWErr．其中，分别计算传感器数量为 100,200,300 的三种情况．

%    如图 14 所示，AREM 的重建误差最低，并且随着传感器数量的增加，误差进一步减小，表明 AREM 能够有效适应高度动态的环境．相比之下，其他方法的重建误差随路径损耗指数的增加而迅速增大．

    \item \textbf{网格分辨率的影响：}
    为验证网格分辨率的影响，假设 ROI 大小为 $800 \times 800\,\mathrm{m}$，并按 10、20、25、40、50 和 $80\,\mathrm{m}$ 的网格间距进行离散化．采样比定义为 $r = M / N$．在采样率固定为 10\% 的情况下，计算对应的 AWErr．

%    我们进行仿真，分析 AREM、VBEM、OMP 和 TPS 的 REM 重建误差与网格间距之间的关系．如图 15 所示，所有方法的重建误差均随网格间距的增大而增加．值得注意的是，OMP 在较大网格间距下对离网效应高度敏感，导致其误差增长速度快于 VBEM 和 AREM．相比之下，考虑了离网效应的 AREM 和 VBEM 实现了显著优于 OMP 的重建精度，其中 AREM 的重建误差最低．
\end{enumerate}

\end{document}